% GENERATED FILE. Edit canonical lessons, then run npm run generate:books.
% canonical-source-hash: fnv1a64:d1059d7df4b3c28d
% canonical-lessons: TA-C483-pun, TA-C483-koppulam, TA-C483-paru, TA-C483-cori, TA-C483-maccam

\chapter{Skin Troubles}
\label{ch:ta-skin-troubles}
\hlchaptermodality{483}

\begin{chapteropening}
\textbf{By the end of this chapter:} \emph{I can talk about sores, blisters, pimples, itching and moles.}

Complete the last lesson of chapter 483: I can talk about sores, blisters, pimples, itching and moles.
\end{chapteropening}

\section[\texorpdfstring{\ta{புண்} (puṇ)}{புண் (puṇ)}]{\ta{புண்} (puṇ) — a sore; an ulcer}

\label{lesson:TA-C483-pun}



\begin{quote}
Before the new one: say the Tamil for to snap, then the Tamil for to ooze.
\end{quote}

\subsection*{You'll want to know: \ta{புண்}}
\textbf{\ta{புண்}} — \emph{puṇ} — \textquotedblleft{}a sore; an ulcer\textquotedblright{}.

\begin{quote}
\ta{வாயில்} \ta{புண்} \ta{இருந்ததால்} \ta{காரமாகச்} \ta{சாப்பிட} \ta{முடியவில்லை}.
\end{quote}

\begin{quote}
\emph{vāyil puṇ iruntatāl kāramākac cāppiṭa muṭiyavillai.}
\end{quote}

\begin{quote}
I had a mouth ulcer, so I could not eat spicy food.
\end{quote}

\ta{வாய்ப்} \ta{புண்}, a mouth ulcer.

\begin{cousinweb}
\textbf{\ta{புண்}} (\emph{puṇ}) is native Dravidian (DEDR 4268), for a wound or sore.
\end{cousinweb}

\begin{grammarlens}[title={What you've built}]
One more word for describing what you lived and felt.
\end{grammarlens}

\subsection*{Your turn}
\begin{itemize}
\raggedright
  \item \emph{Say it:} \emph{puṇ}
  \item \emph{Say it:} \emph{puṇ}, once more
  \item \emph{Say it:} \emph{vāyil puṇ iruntatāl kāramākac cāppiṭa muṭiyavillai.} — I had a mouth ulcer, so I could not eat spicy food.
  \item \emph{Recall:} say the Tamil for to snap, then the Tamil for to ooze, then say \emph{puṇ} again
\end{itemize}

\begin{tcolorbox}[breakable,colback=teal!4,colframe=teal!35!black,title={Before you move on}]
What does \ta{புண்} mean? (\textquotedblleft{}A sore; an ulcer\textquotedblright{}.) Say it once more.
\end{tcolorbox}
\section[\texorpdfstring{\ta{கொப்புளம்} (koppuḷam)}{கொப்புளம் (koppuḷam)}]{\ta{கொப்புளம்} (koppuḷam) — a blister}

\label{lesson:TA-C483-koppulam}



\begin{quote}
Before the new one: say the Tamil for to ooze, then the Tamil for a sore.
\end{quote}

\subsection*{You'll want to know: \ta{கொப்புளம்}}
\textbf{\ta{கொப்புளம்}} — \emph{koppuḷam} — \textquotedblleft{}a blister\textquotedblright{}.

\begin{quote}
\ta{சூடான} \ta{பாத்திரம்} \ta{பட்டுக்} \ta{கையில்} \ta{கொப்புளம்} \ta{வந்தது}.
\end{quote}

\begin{quote}
\emph{cūṭāṉa pāttiram paṭṭuk kaiyil koppuḷam vantatu.}
\end{quote}

\begin{quote}
A hot pot touched my hand and it blistered.
\end{quote}

Its first sense is a bubble.

\begin{cousinweb}
\textbf{\ta{கொப்புளம்}} (\emph{koppuḷam}) is native Dravidian (DEDR 2106): a bubble on the skin.
\end{cousinweb}

\begin{grammarlens}[title={What you've built}]
One more word for describing what you lived and felt.
\end{grammarlens}

\subsection*{Your turn}
\begin{itemize}
\raggedright
  \item \emph{Say it:} \emph{koppuḷam}
  \item \emph{Say it:} \emph{koppuḷam}, once more
  \item \emph{Say it:} \emph{cūṭāṉa pāttiram paṭṭuk kaiyil koppuḷam vantatu.} — A hot pot touched my hand and it blistered.
  \item \emph{Recall:} say the Tamil for to ooze, then the Tamil for a sore, then say \emph{koppuḷam} again
\end{itemize}

\begin{tcolorbox}[breakable,colback=teal!4,colframe=teal!35!black,title={Before you move on}]
What does \ta{கொப்புளம்} mean? (\textquotedblleft{}A blister\textquotedblright{}.) Say it once more.
\end{tcolorbox}
\section[\texorpdfstring{\ta{பரு} (paru)}{பரு (paru)}]{\ta{பரு} (paru) — a pimple}

\label{lesson:TA-C483-paru}



\begin{quote}
Before the new one: say the Tamil for a sore, then the Tamil for a blister.
\end{quote}

\subsection*{You'll want to know: \ta{பரு}}
\textbf{\ta{பரு}} — \emph{paru} — \textquotedblleft{}a pimple\textquotedblright{}.

\begin{quote}
\ta{தேர்வு} \ta{நேரத்தில்} \ta{என்} \ta{முகத்தில்} \ta{பருக்கள்} \ta{வந்தன}.
\end{quote}

\begin{quote}
\emph{tērvu nērattil eṉ mukattil parukkaḷ vantaṉa.}
\end{quote}

\begin{quote}
At exam time I got pimples on my face.
\end{quote}

Plural: \ta{பருக்கள்}, \emph{parukkaḷ}.

Plural \ta{பருக்கள்}, with the doubled \ta{க்}.

\begin{cousinweb}
\textbf{\ta{பரு}} (\emph{paru}) is native Dravidian (DEDR 3974). DEDR keeps \ta{பரு} \textquotedblleft{}to become large\textquotedblright{} as another entry (3972).
\end{cousinweb}

\begin{grammarlens}[title={What you've built}]
One more word for describing what you lived and felt.
\end{grammarlens}

\subsection*{Your turn}
\begin{itemize}
\raggedright
  \item \emph{Say it:} \emph{paru}
  \item \emph{Say it:} \emph{paru}, once more
  \item \emph{Say it:} \emph{tērvu nērattil eṉ mukattil parukkaḷ vantaṉa.} — At exam time I got pimples on my face.
  \item \emph{Recall:} say the Tamil for a sore, then the Tamil for a blister, then say \emph{paru} again
\end{itemize}

\begin{tcolorbox}[breakable,colback=teal!4,colframe=teal!35!black,title={Before you move on}]
What does \ta{பரு} mean? (\textquotedblleft{}A pimple\textquotedblright{}.) Say it once more.
\end{tcolorbox}
\section[\texorpdfstring{\ta{சொறி} (coṟi)}{சொறி (coṟi)}]{\ta{சொறி} (coṟi) — to scratch an itch}

\label{lesson:TA-C483-cori}



\begin{quote}
Before the new one: say the Tamil for a blister, then the Tamil for a pimple.
\end{quote}

\subsection*{You'll want to know: \ta{சொறி}}
\textbf{\ta{சொறி}} — \emph{coṟi} — \textquotedblleft{}to scratch an itch\textquotedblright{}.

\begin{quote}
\ta{கொசு} \ta{கடித்த} \ta{இடத்தை} \ta{அவன்} \ta{சொறிந்தான்}.
\end{quote}

\begin{quote}
\emph{kocu kaṭitta iṭattai avaṉ coṟintāṉ.}
\end{quote}

\begin{quote}
He scratched where the mosquito had bitten.
\end{quote}

Past (I …): \ta{சொறிந்தேன்}, \emph{coṟintēṉ}.

As a noun, an itchy rash.

\begin{cousinweb}
\textbf{\ta{சொறி}} (\emph{coṟi}) is native Dravidian (DEDR 2865), both the itch and the scratching.
\end{cousinweb}

\begin{grammarlens}[title={What you've built}]
One more word for describing what you lived and felt.
\end{grammarlens}

\subsection*{Your turn}
\begin{itemize}
\raggedright
  \item \emph{Say it:} \emph{coṟi}
  \item \emph{Say it:} \emph{coṟi}, once more
  \item \emph{Say it:} \emph{kocu kaṭitta iṭattai avaṉ coṟintāṉ.} — He scratched where the mosquito had bitten.
  \item \emph{Recall:} say the Tamil for a blister, then the Tamil for a pimple, then say \emph{coṟi} again
\end{itemize}

\begin{tcolorbox}[breakable,colback=teal!4,colframe=teal!35!black,title={Before you move on}]
What does \ta{சொறி} mean? (\textquotedblleft{}To scratch an itch\textquotedblright{}.) Say it once more.
\end{tcolorbox}
\section[\texorpdfstring{\ta{மச்சம்} (maccam)}{மச்சம் (maccam)}]{\ta{மச்சம்} (maccam) — a mole}

\label{lesson:TA-C483-maccam}



\begin{quote}
Before the new one: say the Tamil for a pimple, then the Tamil for to scratch an itch.
\end{quote}

\subsection*{You'll want to know: \ta{மச்சம்}}
\textbf{\ta{மச்சம்}} — \emph{maccam} — \textquotedblleft{}a mole\textquotedblright{}.

\begin{quote}
\ta{அவளுடைய} \ta{கன்னத்தில்} \ta{ஒரு} \ta{சிறிய} \ta{மச்சம்} \ta{இருக்கிறது}.
\end{quote}

\begin{quote}
\emph{avaḷuṭaiya kaṉṉattil oru ciṟiya maccam irukkiṟatu.}
\end{quote}

\begin{quote}
She has a small mole on her cheek.
\end{quote}

A mark on the skin, often thought lucky.

\begin{cousinweb}
\textbf{\ta{மச்சம்}} (\emph{maccam}) is native Dravidian (DEDR 4632), in an entry of words for marks and blemishes.
\end{cousinweb}

\begin{grammarlens}[title={What you've built}]
One more word for describing what you lived and felt.
\end{grammarlens}

\subsection*{Your turn}
\begin{itemize}
\raggedright
  \item \emph{Say it:} \emph{maccam}
  \item \emph{Say it:} \emph{maccam}, once more
  \item \emph{Say it:} \emph{avaḷuṭaiya kaṉṉattil oru ciṟiya maccam irukkiṟatu.} — She has a small mole on her cheek.
  \item \emph{Recall:} say the Tamil for a pimple, then the Tamil for to scratch an itch, then say \emph{maccam} again
\end{itemize}

\begin{tcolorbox}[breakable,colback=teal!4,colframe=teal!35!black,title={Before you move on}]
What does \ta{மச்சம்} mean? (\textquotedblleft{}A mole\textquotedblright{}.) Say it once more.
\end{tcolorbox}
